\documentclass[nocopyright]{edoxarticle}
\usepackage{edoxflap}

\begin{document}

\begin{flap}{Lineare Funktionen}{10.09.--24.10.26}
\sprintplanning{Wir können lineare Zusammenhänge darstellen, Geradengleichungen bestimmen und damit Probleme lösen.}

  \stations{%
\stationrow{1}{Koordinatensystem und Punkte}
\stationrow{2}{Steigung am Steigungsdreieck}
\stationrow{3}{Steigung aus zwei Punkten}
\stationrow{4}{Geradengleichung $y=mx+b$}
\stationrow{5}{Gerade aus zwei Punkten}
\stationrow{6}{Schnittpunkt zweier Geraden}
\stationrow{7}{Lineare Modelle / Sachaufgaben}
\stationrow{8}{Individuell: \underline{\hspace{2.5cm}}}%
}

\sprintretro{Alle kommen voran. Wir helfen einander, bis alle die Kernkompetenzen erreichen.}

\parentsign

\end{flap}

\end{document}
